\section*{Capítulo \ref{ch:b3:electrons-in-solids} --- Elétrons nos sólidos}

\begin{solution}{exo:b3:electrons-in-solids:1}
(a) $\tau = m_{\text{e}}/ne^2\rho = \qty{2.5e-14}{s}$. (b) $\ell =
v_{\text{F}}\tau \approx \qty{39}{nm}$. (c) Cerca de 110 \omterm{def:b3:crystalline-solids:lattice}{constantes
de rede}. (d) Um elétron que navega por uma centena de íons sem
espalhar não pode estar ricocheteando nos próprios íons --- a
\omterm{def:b3:crystalline-solids:lattice}{rede} tem de ser transparente para ele, exatamente o que o teorema de
Bloch (\cref{thm:b3:electrons-in-solids:bloch}) prova mais tarde.
\end{solution}

\begin{solution}{exo:b3:electrons-in-solids:2}
(a) $v_{\text{d}} = I/neA = \qty{7.3e-4}{m/s}$ --- menos de um
milímetro por segundo. (b) Cerca de 23 minutos. (c) O \emph{campo}
se estabelece ao longo do fio quase à velocidade da luz e põe
todo o mar de elétrons a derivar de uma vez: os marchadores são lentos,
a ordem de marchar não é. (d) $v_{\text{d}}/v_{\text{F}} \sim
10^{-9}$: a condução é um viés imperceptível sobre um movimento
quântico violento, não um escoamento clássico calmo.
\end{solution}

\begin{solution}{exo:b3:electrons-in-solids:3}
(a) Banda semipreenchida: \omterm{def:b3:electrons-in-solids:classes}{metal}, $\dd\rho/\dd T > 0$. (b) Dois
elétrons preencheriam exatamente a banda dele --- o magnésio conduz
porque a banda cheia dele \emph{se sobrepõe} à próxima, vazia: \omterm{def:b3:electrons-in-solids:classes}{metal},
$\dd\rho/\dd T > 0$. (c) Banda cheia, lacuna de \qty{5.5}{eV}: \omterm{def:b3:electrons-in-solids:classes}{isolante}
(formalmente $\dd\rho/\dd T < 0$, mas essencialmente sem portadores
para contar). (d) \omterm{def:b3:electrons-in-solids:classes}{Semicondutor}: $\dd\rho/\dd T < 0$, a
troca de sinal exponencial que denuncia a classe.
\end{solution}

\begin{solution}{exo:b3:electrons-in-solids:4}
(a) $\lambda = hc/E_{\text{g}} = \qty{1.1}{\micro m}$: o silício é
transparente no infravermelho além disso --- e é por isso que as câmeras
infravermelhas podem ter lentes de silício. (b) \qty{225}{nm}, no
ultravioleta: fóton visível algum pode ser absorvido, e por isso o
diamante é claro como água. (c) \qty{365}{nm}: uma lacuna larga o
bastante para o azul (\qty{2.8}{eV}) existe nos nitretos e em quase
mais nada prático --- o LED azul esperou décadas pelo crescimento de
materiais. (d) A lacuna dele fixa o fóton dele: sobrecarregar acrescenta
calor, não largura de lacuna, e o diodo cozinha no escuro --- a cor de um
LED é uma constante do material, não um ajuste de brilho.
\end{solution}

\begin{solution}{exo:b3:electrons-in-solids:5}
(a) $\kappa/\sigma T = 400/(5.9\times10^7\times300) =
\qty{2.3e-8}{W.\ohm/K^2}$. (b) A dez por cento de $L$. (c)
Os mesmos elétrons da superfície de Fermi carregam a carga e o
calor, e assim o tempo de espalhamento deles se cancela na razão. (d) A
panela de aço conduz calor à comida porque os elétrons dela se movem;
o cabo de madeira, \omterm{def:b3:electrons-in-solids:classes}{isolante} nos dois sentidos, mantém esses elétrons ---
e o calor --- longe da sua mão.
\end{solution}

\begin{solution}{exo:b3:electrons-in-solids:6}
(a) Cinquenta por cento a mais --- flagrantemente ausente do experimento.
(b) $T/T_{\text{F}} \approx 0.004$: o termo eletrônico fica
estrangulado abaixo de um por cento. (c) O $T^3$ da \omterm{def:b3:crystalline-solids:lattice}{rede} desaba
mais depressa que o $T$ dos elétrons: abaixo de alguns kelvins, os
elétrons ``desprezíveis'' são tudo o que resta. (d) $C/T =
\gamma + \beta T^2$: o intercepto $\gamma$ pesa os elétrons,
e a inclinação $\beta$, a \omterm{def:b3:crystalline-solids:lattice}{rede} --- um gráfico, os dois inquilinos.
\end{solution}

\begin{solution}{exo:b3:electrons-in-solids:7}
(a) $k_{\text{F}} = (3\pi^2n)^{1/3} = \qty{1.36e10}{m^{-1}}$:
$\lambda_{\text{F}} = \qty{4.6}{\angstrom}$. (b) $2a =
\qty{7.2}{\angstrom}$: a mesma ordem --- a superfície de Fermi do cobre
chega perto da fronteira da zona, e a física que abre a lacuna
acontece \emph{nas} energias que importam. (c) Uma onda estacionária
acumula a carga dela sobre os íons positivos (energia eletrostática
menor), e a outra, entre eles (maior): mesmo comprimento de onda, duas
energias --- a lacuna. (d) O caminho livre de 110 \omterm{def:b3:crystalline-solids:lattice}{constantes de rede}:
as ondas de Bloch não espalham numa \omterm{def:b3:crystalline-solids:lattice}{rede} perfeita, e assim só
sobram vibrações e impurezas para resistir.
\end{solution}

\begin{solution}{exo:b3:electrons-in-solids:8}
(a) $E_{\text{g}}/2k_{\text{B}} = \qty{6500}{K}$:
$\exp[6500(1/300 - 1/400)] \approx 230$. (b) A resistência do
termistor despenca exponencialmente --- uma curva íngreme e sensível;
a da platina sobe suave e linearmente (mais fônons). Um é um
termômetro por contagem de portadores; o outro, por espalhamento. (c)
Os portadores nascem em \emph{pares} elétron--lacuna, e o
equilíbrio $np = n_{\text{i}}^2$ reparte o custo da lacuna entre
os dois --- daí o $E_{\text{g}}/2$. (d)
O $n_{\text{i}}$ chega a $\qty{5e21}{m^{-3}}$ perto de \qty{760}{K}:
a \omterm{def:b3:electrons-in-solids:doping}{dopagem} se afoga e o circuito esquece o projeto dele. Os dispositivos
reais desistem antes --- as fugas reversas deles, que dobram a cada
dez kelvins, se comportam mal primeiro.
\end{solution}

\begin{solution}{exo:b3:electrons-in-solids:9}
(a) $\qty{5e22}{m^{-3}}$ --- cinco milhões de vezes o $n_{\text{i}}$.
(b) O mesmo fator $\sim5\times10^6$: um grãozinho por milhão de
átomos manda sozinho na condutividade. (c) Porque uma ppm acidental
faria o mesmo sem convite: só um \omterm{def:b3:crystalline-solids:lattice}{cristal} puro a ppb tem uma
condutividade que pertence ao projetista. (d) $d \sim n^{-1/3}
\approx \qty{27}{nm}$, dez vezes a órbita de \qty{2.4}{nm}: os \omterm{def:b3:electrons-in-solids:doping}{doadores}
são hidrogênios isolados; aumente a \omterm{def:b3:electrons-in-solids:doping}{dopagem} cem vezes e as
órbitas se tocam --- uma banda de impurezas e, no fim, um \omterm{def:b3:electrons-in-solids:classes}{metal}.
\end{solution}

\begin{solution}{exo:b3:electrons-in-solids:10}
(a) $E = 13.6\times0.26/11.7^2 \approx \qty{26}{meV}$ (o
valor medido do fósforo, \qty{45}{meV}, mantém a escala
honesta). (b) Comparável a $k_{\text{B}}T = \qty{26}{meV}$:
essencialmente todos os \omterm{def:b3:electrons-in-solids:doping}{doadores} estão ionizados à temperatura ambiente ---
a premissa da \cref{def:b3:electrons-in-solids:doping}. (c) $a =
\qty{0.53}{\angstrom}\times\varepsilon_{\text{r}}/(m^*/
m_{\text{e}}) \approx \qty{2.4}{nm}$. (d) Uma órbita que cobre
dezenas de células vê o \omterm{def:b3:crystalline-solids:lattice}{cristal} como um meio liso --- que é
precisamente a condição sob a qual um $\varepsilon_{\text{r}}$ de volume
e um $m^*$ médio de banda são legítimos:
a aproximação se certifica sozinha.
\end{solution}

\begin{solution}{exo:b3:electrons-in-solids:11}
(a) $V_{\text{bi}} = 0.0259\ln(10^{12}) \approx \qty{0.71}{V}$.
(b) $eV/k_{\text{B}}T = 19.3$: razão $\eu^{19.3} \approx
2\times10^{8}$ --- o diodo é uma via de mão única com oito
dígitos. (c) $I_0 \propto n_{\text{i}}^2 \propto
\eu^{-E_{\text{g}}/k_{\text{B}}T}$: a mesma exponencial do
\cref{exo:b3:electrons-in-solids:8}, daí a
duplicação da regra de bolso. (d) Os quatro diodos formam um losango: a
cada meio ciclo, o par que estiver diretamente polarizado guia a corrente
pela carga no \emph{mesmo} sentido --- CA na entrada, CC
acidentada na saída.
\end{solution}

\begin{solution}{exo:b3:electrons-in-solids:12}
(a) Estreite a lacuna e mais fótons a vencem, mas cada um entrega
apenas a pequena energia da lacuna; alargue-a e cada fóton paga mais,
mas menos se qualificam: colheita $\times$ tensão tem máximo no meio.
(b) Logo abaixo do ótimo de \qty{1.3}{eV}: o teto ideal do silício
é de $\sim30\,\%$ --- perto o bastante para que a barateza dele vença.
(c) A célula de cima, de lacuna larga, pega o azul em alta tensão e
passa o vermelho para a célula de lacuna estreita embaixo: cada fóton é
colhido perto da energia dele, a termalização encolhe e o teto da
pilha sobe rumo aos quarenta. (d) O calor eleva o $I_0$
exponencialmente, arrastando para baixo a tensão de circuito aberto
$\sim(k_{\text{B}}T/e)\ln(I_{\text{L}}/I_0)$; e a própria lacuna
se estreita um pouco, trocando tensão por corrente que ela não recupera
por inteiro.
\end{solution}

\begin{solution}{pb:b3:electrons-in-solids:1}
\textbf{1.} $\lambda_{\text{máx}} \approx \qty{500}{nm}$:
cerca de \qty{2.5}{eV}.
\textbf{2.} $hc/E_{\text{g}} \approx \qty{1.1}{\micro m}$.
\textbf{3.} Ela atravessa as células (não há estado final em que
absorver) e termina como calor no substrato --- aquecendo o telhado, não
os fios.
\textbf{4.} Os $\qty{1.12}{eV}$ sobrevivem como o par; o excesso de
\qty{1.4}{eV} escoa para os fônons em picossegundos --- mais depressa
do que qualquer circuito poderia interceptar.
\textbf{5.} A transparência abaixo da lacuna ($\sim20\,\%$ da potência) e
a termalização acima dela ($\sim30\,\%$): o espectro é taxado pela
metade antes de a junção ver um único elétron.
\textbf{6.} A absorção é um salto quântico de um só fóton que precisa
de um estado final real; em meia lacuna não há nenhum em que pausar, e
os eventos de dois fótons são desprezíveis nas densidades de fótons da luz solar.
\textbf{7.} Os pares criados na \omterm{prop:b3:electrons-in-solids:junction}{zona de depleção} ou perto dela sentem o
campo de barreira: os elétrons são varridos para o lado n, e as lacunas
para o lado p, e o campo os separa mais depressa do que eles conseguem
se encontrar para recombinar --- a carga se acumula até que um circuito
externo a alivie como corrente.
\textbf{8.} $V_{\text{bi}} = 0.0259\ln(10^{12}) \approx
\qty{0.71}{V}$.
\textbf{9.} $60\times\qty{0.6}{V} = \qty{36}{V}$.
\textbf{10.} $I = 400/36 \approx \qty{11}{A}$.
\textbf{11.} No painel inteiro, a luz acima da lacuna é de
$\sim\qty{1400}{W}$ a $\sim\qty{2}{eV}$ cada: $4\times10^{21}$
fótons por segundo --- mas as sessenta células estão em \emph{série},
e assim os \qty{11}{A} (um fluxo de elétrons $I/e \approx
7\times10^{19}\,\unit{s^{-1}}$) passam por cada célula, cuja
parcela de fótons é $4\times10^{21}/60 \approx
7\times10^{19}$ por segundo: quase todo fóton absorvido rende
um elétron coletado. A eficiência quântica é excelente; as
perdas são energéticas, não numéricas.
\textbf{12.} Depois dos 50\,\% do espectro, a célula paga o
déficit de tensão ($0.6/1.12 \approx 54\,\%$) e alguns pontos de
reflexão e recombinação: $0.5\times0.54 \approx 27\,\%$,
aparados até os 20\,\% nominais.
\textbf{13.} A fuga $I_0$: a tensão de circuito aberto
$(k_{\text{B}}T/e)\ln(I_{\text{L}}/I_0)$ para onde a
fotocorrente e a fuga do diodo se equilibram --- cerca de \qty{0.6}{V}, bem
aquém da lacuna.
\textbf{14.} $\sin30^\circ = 0.5$: metade do valor nominal de meio-dia só
pela geometria --- antes das nuvens.
\textbf{15.} $\Delta T = \qty{40}{K}$: $-16\,\%$, deixando
$\approx\qty{335}{W}$ --- os dias mais ensolarados não são os melhores
dias por watt.
\textbf{16.} O calor infla o $I_0$ exponencialmente, e a
tensão de circuito aberto --- a repreensão de um logaritmo --- cai uma
fração de por cento por kelvin; a lacuna ligeiramente encolhida
termina o serviço.
\textbf{17.} A fiação em série impõe uma corrente comum: a célula
sombreada, que não gera nada, é levada à polarização reversa pelas
cinquenta e nove colegas dela e dissipa a potência delas como ponto quente
--- a série se estrangula até a célula mais fraca.
\textbf{18.} O diodo de desvio dá à corrente um caminho direto
em volta da sub-série sombreada, sacrificando a tensão dele em vez da
saída do painel inteiro.
\textbf{19.} $\qty{400}{W}\times0.15\times\qty{8760}{h} \approx
\qty{530}{kWh}$ por ano.
\textbf{20.} $\approx\qty{130}{euros}$ por ano: retorno em
cerca de $300/130 \approx 2.5$ anos, e depois duas décadas de lucro.
\textbf{21.} $500/530 \approx$ um ano: o painel paga a energia de
fabricação dele mais ou menos tão rápido quanto o preço.
\textbf{22.} O preto não é uma escolha, e sim uma \omterm{thm:b3:electrons-in-solids:bloch}{estrutura de bandas}:
a absorção precisa de um estado vazio uma energia de fóton acima de um
cheio, e abaixo da lacuna o silício simplesmente não tem nenhum a
oferecer --- a tinta não acrescenta estados.
\textbf{23.} Empilhe em cima uma célula de lacuna larga para pegar o
azul em alta tensão e passar o vermelho adiante: o tandem bate o
teto de junção única. O empecilho: a pilha em série tem de casar
correntes (e preços) entre as camadas.
\textbf{24.} Tudo o que encurta o $\tau$ e envenena junções:
defeitos criados por UV e impurezas que se difundem para dentro espalham
e aprisionam portadores, a umidade corrói contatos e os pontos quentes
envelhecem os diodos --- a lenta entropia de um \omterm{def:b3:crystalline-solids:lattice}{cristal} a quem se pede
sentar ao sol por trinta anos.
\textbf{25.} A lacuna colhe metade do Sol; a junção transforma
pares em \qty{36}{V} de corrente ordenada; as séries,
a sombra e o calor do verão cobram o quinhão deles; e a
$\sim\qty{530}{kWh}$ por ano contra um retorno energético de um ano e
monetário de três, o árbitro decide: cubra o telhado.
\end{solution}
